\begin{proof}[Proof of \cref{obj:thm:exact-propensity-preservation-boundary}]
\begin{enumerate}
\item Fix \(v=(p,\alpha,\beta,\gamma)\in\mathcal V\), and use the quantities of \cref{obj:def:sharp-frontier-scales}. By \cref{obj:thm:heavy-tail-comparison,obj:thm:oracle-frontier-resolved},
\[
c_v>0,\qquad C_v>0,\qquad
c^{\mathrm e}_v\ge\frac{1}{64\sqrt3}>0,\qquad
C^{\mathrm e}_v>0.
\]
Consequently,
\[
\frac{c_v}{C^{\mathrm e}_v}>0,
\qquad
\frac{C_v}{c^{\mathrm e}_v}>0.
\]
% lean: CausalSmith.Stat.FinitepHomogeneityDensegamma.exact_propensity_preservation_boundary

\item We establish the identities for the comparison exponent. Admissibility gives
\[
p-1>0,\qquad p>0,\qquad q>0,\qquad
\gamma>0,\qquad D_p>0,
\]
and hence
\[
\frac{2\gamma q}{(2\gamma+q)D_p}>0.
\]
The phase identity and channel selection in \cref{obj:thm:heavy-tail-comparison} give
\[
E_4(p)-E_0(p)
=
\frac{2\gamma q}{(2\gamma+q)D_p}\bigl(F(p)-1\bigr),
\]
\[
F(p)\ge1\ \Longrightarrow\ E(p)=E_0(p),
\qquad
F(p)<1\ \Longrightarrow\ E(p)=E_4(p).
\]

If \(F(p)\ge1\), the positive prefactor implies
\[
E_0(p)-E_4(p)\le0.
\]
Thus, using the definition of \(G(v)\),
\[
G(v)=0=E_0(p)-E(p)
=
\frac{2\gamma q}{(2\gamma+q)D_p}\max\{0,1-F(p)\}.
\]
If \(F(p)<1\), the same identity implies
\[
E_0(p)-E_4(p)
=
\frac{2\gamma q}{(2\gamma+q)D_p}\bigl(1-F(p)\bigr)>0.
\]
Channel selection now gives
\[
G(v)=E_0(p)-E_4(p)=E_0(p)-E(p)
=
\frac{2\gamma q}{(2\gamma+q)D_p}\max\{0,1-F(p)\}.
\]
Together, these cases establish
\[
G(v)=E_0(p)-E(p),
\qquad
G(v)=
\frac{2\gamma q}{(2\gamma+q)D_p}\max\{0,1-F(p)\}.
\]
In particular,
\[
F(p)\ge1\ \Longrightarrow\ G(v)=0.
\]
% lean: CausalSmith.Stat.FinitepHomogeneityDensegamma.preservation_gap_identity

\item For every integer \(n\ge2\), the two frontier conclusions in
\cref{obj:thm:heavy-tail-comparison,obj:thm:oracle-frontier-resolved} supply
\[
c_v\rho_n(v)\le r_n^*(v)\le C_v\rho_n(v),
\qquad
c^{\mathrm e}_v\rho_n^{\mathrm e}(v)
\le r_n^{*,\mathrm e}(v)
\le C^{\mathrm e}_v\rho_n^{\mathrm e}(v).
\]
Since \(n>0\), both power scales are positive:
\[
\rho_n(v)=n^{-E(p)}>0,
\qquad
\rho_n^{\mathrm e}(v)=n^{-E_0(p)}>0.
\]
The lower bounds therefore imply
\[
r_n^*(v)>0,\qquad r_n^{*,\mathrm e}(v)>0.
\]
Dividing the original-radius lower bound by the supplied-radius upper bound, and the original-radius upper bound by the supplied-radius lower bound, yields
\[
\frac{c_v\rho_n(v)}
     {C^{\mathrm e}_v\rho_n^{\mathrm e}(v)}
\le
\frac{r_n^*(v)}{r_n^{*,\mathrm e}(v)}
\le
\frac{C_v\rho_n(v)}
     {c^{\mathrm e}_v\rho_n^{\mathrm e}(v)}.
\]
The scale quotient satisfies
\[
\frac{\rho_n(v)}{\rho_n^{\mathrm e}(v)}
=
\frac{n^{-E(p)}}{n^{-E_0(p)}}
=
n^{E_0(p)-E(p)}
=
n^{G(v)}.
\]
Consequently, for every integer \(n\ge2\),
\[
\frac{c_v}{C^{\mathrm e}_v}n^{G(v)}
\le
\frac{r_n^*(v)}{r_n^{*,\mathrm e}(v)}
\le
\frac{C_v}{c^{\mathrm e}_v}n^{G(v)}.
\]
% lean: CausalSmith.Stat.FinitepHomogeneityDensegamma.preservation_ratio_bounds

\item If \(F(p)=1\), the exponent identity gives \(G(v)=0\). Substituting \(n^0=1\) into the quotient bounds gives
\[
\frac{c_v}{C^{\mathrm e}_v}
\le
\frac{r_n^*(v)}{r_n^{*,\mathrm e}(v)}
\le
\frac{C_v}{c^{\mathrm e}_v}
\qquad(n\ge2).
\]
Both bounding constants are strictly positive by \cref{obj:thm:heavy-tail-comparison,obj:thm:oracle-frontier-resolved}.
% lean: CausalSmith.Stat.FinitepHomogeneityDensegamma.exact_propensity_preservation_boundary

\item Suppose \(F(p)<1\). The exponent identity gives
\[
G(v)=
\frac{2\gamma q}{(2\gamma+q)D_p}\bigl(1-F(p)\bigr)>0.
\]
A positive power of the natural sample size tends to infinity, and multiplication by the positive lower multiplier preserves this limit:
\[
n^{G(v)}\longrightarrow+\infty,
\qquad
\frac{c_v}{C^{\mathrm e}_v}n^{G(v)}
\longrightarrow+\infty.
\]
For all sufficiently large natural numbers \(n\), we have \(n\ge2\), so the quotient lower bound gives
\[
\frac{c_v}{C^{\mathrm e}_v}n^{G(v)}
\le
\frac{r_n^*(v)}{r_n^{*,\mathrm e}(v)}.
\]
Comparison with the diverging lower bound proves
\[
\frac{r_n^*(v)}{r_n^{*,\mathrm e}(v)}
\longrightarrow+\infty.
\]
% lean: CausalSmith.Stat.FinitepHomogeneityDensegamma.preservation_ratio_tendsto

Moreover, channel selection gives \(E(p)=E_4(p)\), and hence
\[
G(v)=E_0(p)-E(p)=E_0(p)-E_4(p).
\]
Substitution into the quotient bounds proves
\[
\frac{c_v}{C^{\mathrm e}_v}n^{E_0(p)-E_4(p)}
\le
\frac{r_n^*(v)}{r_n^{*,\mathrm e}(v)}
\le
\frac{C_v}{c^{\mathrm e}_v}n^{E_0(p)-E_4(p)}
\qquad(n\ge2).
\]
% lean: CausalSmith.Stat.FinitepHomogeneityDensegamma.exact_propensity_preservation_boundary

\item We characterize the preservation region of
\cref{obj:def:preservation-region}. Suppose first that \(v\in\mathcal R\).
There is a real upper bound \(b_{\mathrm{sup}}\) satisfying
\[
\frac{r_n^*(v)}{r_n^{*,\mathrm e}(v)}
\le b_{\mathrm{sup}}
\qquad\text{for every integer }n\ge2.
\]
If \(F(p)<1\), the divergence just proved supplies an integer \(n\ge2\) such that
\[
b_{\mathrm{sup}}
<
\frac{r_n^*(v)}{r_n^{*,\mathrm e}(v)}
\le b_{\mathrm{sup}},
\]
a contradiction. Thus \(F(p)\ge1\).

Conversely, if \(v\in\mathcal V\) and \(F(p)\ge1\), then \(G(v)=0\), so
\[
\frac{r_n^*(v)}{r_n^{*,\mathrm e}(v)}
\le\frac{C_v}{c^{\mathrm e}_v}
\qquad\text{for every integer }n\ge2.
\]
This finite upper bound places \(v\) in \(\mathcal R\). Expanding \(F(p)\) therefore gives
\[
\mathcal R
=
\left\{v\in\mathcal V:
\frac{2\alpha}{p-1}
+\frac{\beta}{q}
+\frac{\alpha+\beta}{2\gamma}\ge1
\right\}.
\]
% lean: CausalSmith.Stat.FinitepHomogeneityDensegamma.exact_propensity_preservation_boundary

\item Finally, let \(\mathcal K\subseteq\mathcal V\) be an arbitrary compact set.
The compact-uniformity conclusion of
\cref{obj:thm:heavy-tail-comparison} supplies real constants
\(\underline c_{\mathcal K},\overline C_{\mathcal K}\) satisfying
\[
\underline c_{\mathcal K}>0,
\qquad
\underline c_{\mathcal K}\le c_v,
\qquad
C_v\le\overline C_{\mathcal K}
\qquad(v\in\mathcal K).
\]
The supplied-propensity upper multiplier is independent of \(v\) and positive. Define
\[
c_{\mathrm{pres},\mathcal K}
=
\frac{\underline c_{\mathcal K}}
{3\bigl(120+128\sqrt{160\sqrt2}\bigr)},
\qquad
C_{\mathrm{pres},\mathcal K}
=
64\sqrt3\,\overline C_{\mathcal K}.
\]
Then \(c_{\mathrm{pres},\mathcal K}>0\), and for every \(v\in\mathcal K\),
\[
c_{\mathrm{pres},\mathcal K}
\le\frac{c_v}{C^{\mathrm e}_v},
\qquad
64\sqrt3\,C_v
\le C_{\mathrm{pres},\mathcal K}.
\]
These are the required compact-uniform constants.
% lean: CausalSmith.Stat.FinitepHomogeneityDensegamma.exact_propensity_preservation_boundary
\end{enumerate}
\end{proof}
